\makeatletter
\newlabel{fig:haz}{{1.1}{2}{Búsqueda en haz con $m=2$ en el ejemplo de la \cref {prop:voraz-no-optima}. Los nodos discontinuos se podan (no se dibujan las extensiones con $\mathrm {z}$). La búsqueda voraz se queda con $A\mathrm {x}$ ($0{,}20$); el haz encuentra $B\mathrm {x}$ ($0{,}36$)}{figure.1.1}{}}
\newlabel{fig:haz@cref}{{[figure][1][1]1.1}{[1][2][]2}}
\newlabel{ej:muestreo}{{1.12}{3}{}{ejemplo.12}{}}
\newlabel{ej:muestreo@cref}{{[ejemplo][12][]1.12}{[1][3][]3}}
\newlabel{ej:truncamiento}{{1.17}{4}{}{ejemplo.17}{}}
\newlabel{ej:truncamiento@cref}{{[ejemplo][17][]1.17}{[1][4][]4}}
\newlabel{fig:cadena-muestreo}{{1.2}{5}{Cadena habitual de muestreo. Con $\tau \to 0$ o con $k=1$ se recupera la búsqueda voraz}{figure.1.2}{}}
\newlabel{fig:cadena-muestreo@cref}{{[figure][2][1]1.2}{[1][5][]5}}
\newlabel{eje:1-topk-topp}{{1.1}{5}{}{ejercicio.1.1}{}}
\newlabel{eje:1-topk-topp@cref}{{[ejercicio][1][1]1.1}{[1][5][]5}}
\newlabel{eje:1-logits}{{1.2}{5}{}{ejercicio.1.2}{}}
\newlabel{eje:1-logits@cref}{{[ejercicio][2][1]1.2}{[1][5][]5}}
\newlabel{eje:1-slp}{{1.3}{5}{}{ejercicio.1.3}{}}
\newlabel{eje:1-slp@cref}{{[ejercicio][3][1]1.3}{[1][5][]5}}
\newlabel{eje:1-nucleo}{{1.4}{5}{}{ejercicio.1.4}{}}
\newlabel{eje:1-nucleo@cref}{{[ejercicio][4][1]1.4}{[1][5][]5}}
\newlabel{eje:1-haz}{{1.5}{5}{}{ejercicio.1.5}{}}
\newlabel{eje:1-haz@cref}{{[ejercicio][5][1]1.5}{[1][5][]5}}
\newlabel{eje:1-potencia}{{1.6}{6}{}{ejercicio.1.6}{}}
\newlabel{eje:1-potencia@cref}{{[ejercicio][6][1]1.6}{[1][6][]6}}
\newlabel{eje:1-entropia}{{1.7}{6}{}{ejercicio.1.7}{}}
\newlabel{eje:1-entropia@cref}{{[ejercicio][7][1]1.7}{[1][6][]6}}
\newlabel{ej:prefill}{{2.2}{7}{Coste del \emph {prefill}}{ejemplo.2}{}}
\newlabel{ej:prefill@cref}{{[ejemplo][2][]2.2}{[1][7][]7}}
\newlabel{fig:cache}{{2.1}{8}{Caché KV de una capa. El \emph {prefill} rellena las $n$ primeras columnas; cada paso de decodificación calcula solo la consulta, la clave y el valor de la posición nueva, añade una columna y atiende sobre todas}{figure.2.1}{}}
\newlabel{fig:cache@cref}{{[figure][1][2]2.1}{[1][8][]8}}
\newlabel{ej:llama}{{2.7}{9}{Arquitecturas públicas}{ejemplo.7}{}}
\newlabel{ej:llama@cref}{{[ejemplo][7][]2.7}{[1][9][]9}}
\newlabel{ej:cota}{{2.11}{10}{Cota de velocidad}{ejemplo.11}{}}
\newlabel{ej:cota@cref}{{[ejemplo][11][]2.11}{[1][10][]10}}
\newlabel{eje:2-fases}{{2.1}{11}{}{ejercicio.2.1}{}}
\newlabel{eje:2-fases@cref}{{[ejercicio][1][2]2.1}{[1][11][]11}}
\newlabel{eje:2-vf}{{2.2}{11}{}{ejercicio.2.2}{}}
\newlabel{eje:2-vf@cref}{{[ejercicio][2][2]2.2}{[1][11][]11}}
\newlabel{eje:2-kv}{{2.3}{11}{}{ejercicio.2.3}{}}
\newlabel{eje:2-kv@cref}{{[ejercicio][3][2]2.3}{[1][11][]11}}
\newlabel{eje:2-cota}{{2.4}{11}{}{ejercicio.2.4}{}}
\newlabel{eje:2-cota@cref}{{[ejercicio][4][2]2.4}{[1][11][]11}}
\newlabel{eje:2-llama13}{{2.5}{11}{}{ejercicio.2.5}{}}
\newlabel{eje:2-llama13@cref}{{[ejercicio][5][2]2.5}{[1][11][]11}}
\newlabel{eje:2-intensidad}{{2.6}{11}{}{ejercicio.2.6}{}}
\newlabel{eje:2-intensidad@cref}{{[ejercicio][6][2]2.6}{[1][11][]11}}
\newlabel{eje:2-sincache}{{2.7}{12}{}{ejercicio.2.7}{}}
\newlabel{eje:2-sincache@cref}{{[ejercicio][7][2]2.7}{[1][11][]12}}
\newlabel{fig:ronda}{{3.1}{14}{Una ronda con $\gamma =4$. El borrador propone cuatro tokens (arriba); el objetivo los evalúa en una pasada y acepta los dos primeros. En el primer rechazo emite un token de la distribución residual y la ronda termina}{figure.3.1}{}}
\newlabel{fig:ronda@cref}{{[figure][1][3]3.1}{[1][13][]14}}
\newlabel{ej:especulativa}{{3.6}{15}{}{ejemplo.6}{}}
\newlabel{ej:especulativa@cref}{{[ejemplo][6][]3.6}{[1][14][]15}}
\newlabel{ej:gamma}{{3.13}{16}{Elección de $\gamma $}{ejemplo.13}{}}
\newlabel{ej:gamma@cref}{{[ejemplo][13][]3.13}{[1][16][]16}}
\newlabel{eje:3-vf}{{3.1}{16}{}{ejercicio.3.1}{}}
\newlabel{eje:3-vf@cref}{{[ejercicio][1][3]3.1}{[1][16][]16}}
\newlabel{eje:3-dos}{{3.2}{16}{}{ejercicio.3.2}{}}
\newlabel{eje:3-dos@cref}{{[ejercicio][2][3]3.2}{[1][16][]16}}
\newlabel{eje:3-cs336}{{3.3}{17}{}{ejercicio.3.3}{}}
\newlabel{eje:3-cs336@cref}{{[ejercicio][3][3]3.3}{[1][17][]17}}
\newlabel{eje:3-gamma}{{3.4}{17}{}{ejercicio.3.4}{}}
\newlabel{eje:3-gamma@cref}{{[ejercicio][4][3]3.4}{[1][17][]17}}
\newlabel{eje:3-tv}{{3.5}{17}{}{ejercicio.3.5}{}}
\newlabel{eje:3-tv@cref}{{[ejercicio][5][3]3.5}{[1][17][]17}}
\newlabel{eje:3-optimo}{{3.6}{17}{}{ejercicio.3.6}{}}
\newlabel{eje:3-optimo@cref}{{[ejercicio][6][3]3.6}{[1][17][]17}}
\newlabel{fig:rejilla}{{4.1}{19}{Rejilla del cuantizador afín de 4 bits del \cref {eje:4-int4}: los 16 enteros $Q=0,\dots ,15$ cubren $[-1,2]$ con paso $\Delta =0{,}2$; el real $0$ corresponde al punto cero $Z=5$. El valor $x=0{,}62$ cae en $x/\Delta +Z=8{,}1$ y se lleva al entero $Q=8$, que se reconstruye como $\hat x=0{,}6$}{figure.4.1}{}}
\newlabel{fig:rejilla@cref}{{[figure][1][4]4.1}{[1][18][]19}}
\newlabel{ej:cuant}{{4.4}{19}{}{ejemplo.4}{}}
\newlabel{ej:cuant@cref}{{[ejemplo][4][]4.4}{[1][19][]19}}
\newlabel{eje:4-conceptos}{{4.1}{21}{}{ejercicio.4.1}{}}
\newlabel{eje:4-conceptos@cref}{{[ejercicio][1][4]4.1}{[1][21][]21}}
\newlabel{eje:4-int4}{{4.2}{21}{}{ejercicio.4.2}{}}
\newlabel{eje:4-int4@cref}{{[ejercicio][2][4]4.2}{[1][21][]21}}
\newlabel{eje:4-3bits}{{4.3}{21}{}{ejercicio.4.3}{}}
\newlabel{eje:4-3bits@cref}{{[ejercicio][3][4]4.3}{[1][21][]21}}
\newlabel{eje:4-atipico}{{4.4}{22}{}{ejercicio.4.4}{}}
\newlabel{eje:4-atipico@cref}{{[ejercicio][4][4]4.4}{[1][22][]22}}
\newlabel{eje:4-entero}{{4.5}{22}{}{ejercicio.4.5}{}}
\newlabel{eje:4-entero@cref}{{[ejercicio][5][4]4.5}{[1][22][]22}}
\newlabel{eje:4-ruido}{{4.6}{22}{}{ejercicio.4.6}{}}
\newlabel{eje:4-ruido@cref}{{[ejercicio][6][4]4.6}{[1][22][]22}}
\newlabel{eje:4-granularidad}{{4.7}{22}{}{ejercicio.4.7}{}}
\newlabel{eje:4-granularidad@cref}{{[ejercicio][7][4]4.7}{[1][22][]22}}
\newlabel{eje:5-conceptos}{{5.1}{27}{}{ejercicio.5.1}{}}
\newlabel{eje:5-conceptos@cref}{{[ejercicio][1][5]5.1}{[1][27][]27}}
\newlabel{eje:5-absmax}{{5.2}{27}{}{ejercicio.5.2}{}}
\newlabel{eje:5-absmax@cref}{{[ejercicio][2][5]5.2}{[1][27][]27}}
\newlabel{eje:5-smooth}{{5.3}{27}{}{ejercicio.5.3}{}}
\newlabel{eje:5-smooth@cref}{{[ejercicio][3][5]5.3}{[1][27][]27}}
\newlabel{eje:5-gptq}{{5.4}{28}{}{ejercicio.5.4}{}}
\newlabel{eje:5-gptq@cref}{{[ejercicio][4][5]5.4}{[1][28][]28}}
\newlabel{eje:5-qlora}{{5.5}{28}{}{ejercicio.5.5}{}}
\newlabel{eje:5-qlora@cref}{{[ejercicio][5][5]5.5}{[1][28][]28}}
\newlabel{eje:5-destila}{{5.6}{28}{}{ejercicio.5.6}{}}
\newlabel{eje:5-destila@cref}{{[ejercicio][6][5]5.6}{[1][28][]28}}
\newlabel{eje:5-altaT}{{5.7}{28}{}{ejercicio.5.7}{}}
\newlabel{eje:5-altaT@cref}{{[ejercicio][7][5]5.7}{[1][28][]28}}
\newlabel{ej:latencia}{{6.3}{29}{}{ejemplo.3}{}}
\newlabel{ej:latencia@cref}{{[ejemplo][3][]6.3}{[1][29][]29}}
\newlabel{fig:batching}{{6.1}{30}{Tres ranuras durante ocho pasos de decodificación. Con \emph {batching} estático (izquierda), las ranuras de $A$ y $C$ quedan vacías (rayado) hasta que termina $B$: se usan 16 de 24 ranuras. Con \emph {batching} continuo (derecha), $D$ y $E$ entran en cuanto se liberan las ranuras y se usan las 24}{figure.6.1}{}}
\newlabel{fig:batching@cref}{{[figure][1][6]6.1}{[1][30][]30}}
\newlabel{fig:paged}{{6.2}{31}{Una secuencia de 40 tokens con bloques de 16. La tabla de bloques ($0\mapsto 5$, $1\mapsto 1$, $2\mapsto 6$) la sitúa en tres bloques físicos no contiguos; solo se desperdician 8 posiciones del último bloque}{figure.6.2}{}}
\newlabel{fig:paged@cref}{{[figure][2][6]6.2}{[1][31][]31}}
\newlabel{ej:bloques}{{6.10}{31}{}{ejemplo.10}{}}
\newlabel{ej:bloques@cref}{{[ejemplo][10][]6.10}{[1][31][]31}}
\newlabel{eje:6-paginacion}{{6.1}{32}{}{ejercicio.6.1}{}}
\newlabel{eje:6-paginacion@cref}{{[ejercicio][1][6]6.1}{[1][32][]32}}
\newlabel{eje:6-bloques}{{6.2}{32}{}{ejercicio.6.2}{}}
\newlabel{eje:6-bloques@cref}{{[ejercicio][2][6]6.2}{[1][32][]32}}
\newlabel{eje:6-latencia}{{6.3}{32}{}{ejercicio.6.3}{}}
\newlabel{eje:6-latencia@cref}{{[ejercicio][3][6]6.3}{[1][32][]32}}
\newlabel{eje:6-gqa}{{6.4}{32}{}{ejercicio.6.4}{}}
\newlabel{eje:6-gqa@cref}{{[ejercicio][4][6]6.4}{[1][32][]32}}
\newlabel{eje:6-prefijo}{{6.5}{32}{}{ejercicio.6.5}{}}
\newlabel{eje:6-prefijo@cref}{{[ejercicio][5][6]6.5}{[1][32][]32}}
\newlabel{eje:6-frag}{{6.6}{32}{}{ejercicio.6.6}{}}
\newlabel{eje:6-frag@cref}{{[ejercicio][6][6]6.6}{[1][32][]32}}
\newlabel{ap:soluciones}{{A}{33}{Soluciones de los ejercicios}{appendix.A}{}}
\newlabel{ap:soluciones@cref}{{[appendix][1][2147483647]A}{[1][33][]33}}
\newlabel{sol:eje:1-topk-topp}{{A.1}{33}{Tema 1}{section*.3}{}}
\newlabel{sol:eje:1-topk-topp@cref}{{[subappendix][1][2147483647,1]A.1}{[1][33][]33}}
\newlabel{sol:eje:1-logits}{{A.1}{33}{Tema 1}{section*.4}{}}
\newlabel{sol:eje:1-logits@cref}{{[subappendix][1][2147483647,1]A.1}{[1][33][]33}}
\newlabel{sol:eje:1-slp}{{A.1}{34}{Tema 1}{section*.5}{}}
\newlabel{sol:eje:1-slp@cref}{{[subappendix][1][2147483647,1]A.1}{[1][34][]34}}
\newlabel{sol:eje:1-nucleo}{{A.1}{34}{Tema 1}{section*.6}{}}
\newlabel{sol:eje:1-nucleo@cref}{{[subappendix][1][2147483647,1]A.1}{[1][34][]34}}
\newlabel{sol:eje:1-haz}{{A.1}{35}{Tema 1}{section*.7}{}}
\newlabel{sol:eje:1-haz@cref}{{[subappendix][1][2147483647,1]A.1}{[1][35][]35}}
\newlabel{sol:eje:1-potencia}{{A.1}{35}{Tema 1}{section*.8}{}}
\newlabel{sol:eje:1-potencia@cref}{{[subappendix][1][2147483647,1]A.1}{[1][35][]35}}
\newlabel{sol:eje:1-entropia}{{A.1}{36}{Tema 1}{section*.9}{}}
\newlabel{sol:eje:1-entropia@cref}{{[subappendix][1][2147483647,1]A.1}{[1][36][]36}}
\newlabel{sol:eje:2-fases}{{A.2}{36}{Tema 2}{section*.10}{}}
\newlabel{sol:eje:2-fases@cref}{{[subappendix][2][2147483647,1]A.2}{[1][36][]36}}
\newlabel{sol:eje:2-vf}{{A.2}{37}{Tema 2}{section*.11}{}}
\newlabel{sol:eje:2-vf@cref}{{[subappendix][2][2147483647,1]A.2}{[1][37][]37}}
\newlabel{sol:eje:2-kv}{{A.2}{37}{Tema 2}{section*.12}{}}
\newlabel{sol:eje:2-kv@cref}{{[subappendix][2][2147483647,1]A.2}{[1][37][]37}}
\newlabel{sol:eje:2-cota}{{A.2}{37}{Tema 2}{section*.13}{}}
\newlabel{sol:eje:2-cota@cref}{{[subappendix][2][2147483647,1]A.2}{[1][37][]37}}
\newlabel{sol:eje:2-llama13}{{A.2}{38}{Tema 2}{section*.14}{}}
\newlabel{sol:eje:2-llama13@cref}{{[subappendix][2][2147483647,1]A.2}{[1][37][]38}}
\newlabel{sol:eje:2-intensidad}{{A.2}{38}{Tema 2}{section*.15}{}}
\newlabel{sol:eje:2-intensidad@cref}{{[subappendix][2][2147483647,1]A.2}{[1][38][]38}}
\newlabel{sol:eje:2-sincache}{{A.2}{39}{Tema 2}{section*.16}{}}
\newlabel{sol:eje:2-sincache@cref}{{[subappendix][2][2147483647,1]A.2}{[1][38][]39}}
\newlabel{sol:eje:3-vf}{{A.3}{39}{Tema 3}{section*.17}{}}
\newlabel{sol:eje:3-vf@cref}{{[subappendix][3][2147483647,1]A.3}{[1][39][]39}}
\newlabel{sol:eje:3-dos}{{A.3}{39}{Tema 3}{section*.18}{}}
\newlabel{sol:eje:3-dos@cref}{{[subappendix][3][2147483647,1]A.3}{[1][39][]39}}
\newlabel{sol:eje:3-cs336}{{A.3}{40}{Tema 3}{section*.19}{}}
\newlabel{sol:eje:3-cs336@cref}{{[subappendix][3][2147483647,1]A.3}{[1][40][]40}}
\newlabel{sol:eje:3-gamma}{{A.3}{40}{Tema 3}{section*.20}{}}
\newlabel{sol:eje:3-gamma@cref}{{[subappendix][3][2147483647,1]A.3}{[1][40][]40}}
\newlabel{sol:eje:3-tv}{{A.3}{40}{Tema 3}{section*.21}{}}
\newlabel{sol:eje:3-tv@cref}{{[subappendix][3][2147483647,1]A.3}{[1][40][]40}}
\newlabel{sol:eje:3-optimo}{{A.3}{41}{Tema 3}{section*.22}{}}
\newlabel{sol:eje:3-optimo@cref}{{[subappendix][3][2147483647,1]A.3}{[1][41][]41}}
\newlabel{sol:eje:4-conceptos}{{A.4}{41}{Tema 4}{section*.23}{}}
\newlabel{sol:eje:4-conceptos@cref}{{[subappendix][4][2147483647,1]A.4}{[1][41][]41}}
\newlabel{sol:eje:4-int4}{{A.4}{42}{Tema 4}{section*.24}{}}
\newlabel{sol:eje:4-int4@cref}{{[subappendix][4][2147483647,1]A.4}{[1][42][]42}}
\newlabel{sol:eje:4-3bits}{{A.4}{42}{Tema 4}{section*.25}{}}
\newlabel{sol:eje:4-3bits@cref}{{[subappendix][4][2147483647,1]A.4}{[1][42][]42}}
\newlabel{sol:eje:4-atipico}{{A.4}{43}{Tema 4}{section*.26}{}}
\newlabel{sol:eje:4-atipico@cref}{{[subappendix][4][2147483647,1]A.4}{[1][43][]43}}
\newlabel{sol:eje:4-entero}{{A.4}{43}{Tema 4}{section*.27}{}}
\newlabel{sol:eje:4-entero@cref}{{[subappendix][4][2147483647,1]A.4}{[1][43][]43}}
\newlabel{sol:eje:4-ruido}{{A.4}{44}{Tema 4}{section*.28}{}}
\newlabel{sol:eje:4-ruido@cref}{{[subappendix][4][2147483647,1]A.4}{[1][44][]44}}
\newlabel{sol:eje:4-granularidad}{{A.4}{44}{Tema 4}{section*.29}{}}
\newlabel{sol:eje:4-granularidad@cref}{{[subappendix][4][2147483647,1]A.4}{[1][44][]44}}
\newlabel{sol:eje:5-conceptos}{{A.5}{45}{Tema 5}{section*.30}{}}
\newlabel{sol:eje:5-conceptos@cref}{{[subappendix][5][2147483647,1]A.5}{[1][44][]45}}
\newlabel{sol:eje:5-absmax}{{A.5}{45}{Tema 5}{section*.31}{}}
\newlabel{sol:eje:5-absmax@cref}{{[subappendix][5][2147483647,1]A.5}{[1][45][]45}}
\newlabel{sol:eje:5-smooth}{{A.5}{45}{Tema 5}{section*.32}{}}
\newlabel{sol:eje:5-smooth@cref}{{[subappendix][5][2147483647,1]A.5}{[1][45][]45}}
\newlabel{sol:eje:5-gptq}{{A.5}{46}{Tema 5}{section*.33}{}}
\newlabel{sol:eje:5-gptq@cref}{{[subappendix][5][2147483647,1]A.5}{[1][46][]46}}
\newlabel{sol:eje:5-qlora}{{A.5}{47}{Tema 5}{section*.34}{}}
\newlabel{sol:eje:5-qlora@cref}{{[subappendix][5][2147483647,1]A.5}{[1][47][]47}}
\newlabel{sol:eje:5-destila}{{A.5}{47}{Tema 5}{section*.35}{}}
\newlabel{sol:eje:5-destila@cref}{{[subappendix][5][2147483647,1]A.5}{[1][47][]47}}
\newlabel{sol:eje:5-altaT}{{A.5}{48}{Tema 5}{section*.36}{}}
\newlabel{sol:eje:5-altaT@cref}{{[subappendix][5][2147483647,1]A.5}{[1][47][]48}}
\newlabel{sol:eje:6-paginacion}{{A.6}{48}{Tema 6}{section*.37}{}}
\newlabel{sol:eje:6-paginacion@cref}{{[subappendix][6][2147483647,1]A.6}{[1][48][]48}}
\newlabel{sol:eje:6-bloques}{{A.6}{49}{Tema 6}{section*.38}{}}
\newlabel{sol:eje:6-bloques@cref}{{[subappendix][6][2147483647,1]A.6}{[1][49][]49}}
\newlabel{sol:eje:6-latencia}{{A.6}{49}{Tema 6}{section*.39}{}}
\newlabel{sol:eje:6-latencia@cref}{{[subappendix][6][2147483647,1]A.6}{[1][49][]49}}
\newlabel{sol:eje:6-gqa}{{A.6}{49}{Tema 6}{section*.40}{}}
\newlabel{sol:eje:6-gqa@cref}{{[subappendix][6][2147483647,1]A.6}{[1][49][]49}}
\newlabel{sol:eje:6-prefijo}{{A.6}{50}{Tema 6}{section*.41}{}}
\newlabel{sol:eje:6-prefijo@cref}{{[subappendix][6][2147483647,1]A.6}{[1][50][]50}}
\newlabel{sol:eje:6-frag}{{A.6}{50}{Tema 6}{section*.42}{}}
\newlabel{sol:eje:6-frag@cref}{{[subappendix][6][2147483647,1]A.6}{[1][50][]50}}
\makeatother
